% font: footnotesize
% pack: with-previous
% TAB-10: generated by harness/analysis/paper_outputs.py; do not edit (every cell's source is in TAB-10.cells.json)
\begin{tabularx}{\linewidth}{l>{\hsize=0.773\hsize\raggedleft\hyphenpenalty=10000\exhyphenpenalty=10000\arraybackslash}X>{\hsize=1.310\hsize\raggedleft\hyphenpenalty=10000\exhyphenpenalty=10000\arraybackslash}X>{\hsize=0.743\hsize\raggedleft\hyphenpenalty=10000\exhyphenpenalty=10000\arraybackslash}X>{\hsize=1.174\hsize\raggedleft\hyphenpenalty=10000\exhyphenpenalty=10000\arraybackslash}X}
\toprule
Test & 10-point margin: $p$ & 10-point margin: Holm threshold & Error-rate ratio $R$ & Error-rate ratio 95\% interval \\
\midrule
Test~B1 & 0.00027 & 0.0125\phantom{0} & 3.37 & \mbox{[1.78, 10.07]} \\
Test~B2 & \ensuremath{\leq}~0.00011 & 0.00833 & 0.71 & \mbox{[0.24, 2.11]} \\
Test~B3 & \ensuremath{\leq}~0.00108 & 0.025\phantom{00} & 6.80 & \mbox{[1.67, unbounded]} \\
\bottomrule
\end{tabularx}
% note: Readings in quotation marks are the Study~1b file's own template words: ``holds'' only if the test is no worse in all three views; an inconclusive test reads ``inconclusive (could not show)'' and is never called worse. The ruled reading, shown no worse (Holm-adjusted), is the column ``5-point margin: shown no worse (Holm-adjusted)'' of this table. Panel~2 has no such column: at the 5-point margin Test~B2 is shown no worse (Holm-adjusted) and Tests~B1 and B3 are not.
% note: Holm's correction runs within Study~1b only; Study~1 and Study~1b are never pooled. Shown no worse (Holm-adjusted) is yes where $p$ is below the test's Holm threshold at that margin, in Holm's step-down order: a test is shown no worse (Holm-adjusted) only if every test before it in Holm's order is too.
% note: $d$ is the same at both margins: Haiku-max (or its vote) minus Opus-medium, in points, paired over puzzles. A positive $d$ means Haiku-max scored higher.
% note: Error-rate ratio $R$ is the challenger's error rate divided by Opus-medium's, over the same puzzles (an error rate is the share of answers not right, the complement of the accuracy); $R$ above one means the challenger made more errors. $R$ is pre-specified descriptive, reported beside the test: Study~1b's specification, frozen with Lock~4, reports it beside every test, and it is not part of any test's reading. Its 95\% interval is a bootstrap over puzzles; an unbounded upper end means Opus-medium made no error in too many resamples for the interval to have a finite top.
% note: The B2 challenger is the vote of $k$ Haiku-max samples, $k$~=~6 (cost-matched with Opus-medium, by the pre-registered cost-matched $k$ rule).
% note: For Test~B2 the challenger's format compliance is that of the single Haiku-max samples the vote is taken over.
% note: Costs: list price per answer; the vote at $k$ costs $k$ answers. Each test's cost wording, $k$ and per-shape costs (B2's included) are in Panel~1 (costs). Its ratio column holds Test~B2's spend share (the vote's cost over its cost-matched comparator's) and the cost ratios of Tests~B1 and B3.
% note: Panel~2 (sensitivity) leaves out the puzzles that were exposed before the run (the file lists them under excluded); it is a sensitivity, not a test.
% note: Panel~1's wordings are the Study~1b file's own, in quotation marks; its ``tries'' are the paper's samples.
% note: Columns that print one value in every row are left out. In every test the comparator is Opus-medium, whose format compliance is 100.00\% / 100.00\% (chains / pieces); at the 10-point margin every test reads ``holds (no worse in~all~three~views)'' and is shown no worse (Holm-adjusted). Panel~2's 10-point readings are the same.
% note: \ensuremath{\leq}~marks a $p$ at the smallest value the bootstrap can return (no resample reached the margin).
